\section{Discussion}
\label{sec:discussion}

Four findings, and each is bounded by something the design cannot recover. Motor fuel is the largest single channel at \genMotorFuelShareOfLoss{} per cent of the aggregate loss, but between \genSpecFuelShareMinPct{} and \genSpecFuelShareMaxPct{} per cent across specifications and wider still across the gas-peak profiles the pre-war solve consumes; what holds across that range is that a large part of the shock, never less than \genSpecFuelShareMinPct{} per cent, arrives through a channel every instrument in the Budget debate leaves untouched. The cap formula can carry a macro scenario to a household price vector, but only against a counterfactual it does not supply: the pre-war level is solved, at \pounds \genCapBaselineLevel{}, and its fragility to the gas-peak profile --- \pounds \genGasProfileAggMinBn{} to \pounds \genGasProfileAggMaxBn{} billion, with \genGasProfileNonIdentifiedCount{} of \genGasProfileCellCount{} profiles admitting no solution --- is the paper's principal uncertainty. The burden falls in income in every specification while the cash loss does not, though the steepness is uncertain by more than a factor of two and every gradient here is ranked on before-housing-costs income and measured on after-housing-costs income (Section~\ref{sec:limitations}). And eligibility rather than generosity is the binding margin on an exchequer envelope, by a factor the means-tested fuel imputation does not pin down.

Two of those results deserve a note about what stabilises them. Measured on the domestic channel alone the decile gradient is \genDomesticOnlyDecileRatioPct{} to one in every one of the seven specifications and every cell of the scenario grid, so the variation in the all-channel figure is located in the fuel imputation rather than in the domestic leg. And the eligibility finding survives every admission rule scored, including a lottery: what the perfect-observability assumption buys is targeting, not offset. No constituency-level claim is made anywhere: the data release has no constituency weight matrix and a degenerate local-authority field, and no approximation is offered in place of one.

\subsection{The K\"anzig caveat}

The most serious objection is that this paper measures the wrong channel. \citet{kanzig2023} shows that poorer households lose from energy shocks disproportionately through employment and earnings rather than budget shares, with general-equilibrium and indirect effects accounting for roughly two-thirds of the household response, none of which a static use-side microsimulation captures. \citet{goulder2019} sharpen the point --- use-side effects regressive, omitted source-side effects progressive --- and \citet{ohlendorf2021}'s meta-analysis is the summary statement. It constrains the claims here in three ways.

\emph{Magnitude.} If the use-side channel is roughly a third of the total, the losses reported are a lower bound on total household incidence, and a bottom-decile loss of \genDecileOneLossPct{} per cent of income is the part attributable to the price of energy alone. A reader wanting a total should multiply, acknowledging that the multiplier is borrowed from a different country and identification strategy.

\emph{Ordering.} Känzig's mechanism operates through employment, which is more concentrated among lower-income households than energy budget shares are, so the omitted channel almost certainly reinforces the percentage regressivity --- but it may well \emph{reverse} the cash ordering, since employment losses in pounds are larger higher up the earnings distribution. A general-equilibrium version of this exercise would not rescale the cash profile; it would redraw it.

\emph{Policy inference.} The compensation metrics are internally consistent, so the ranking of policies is more robust than the levels, provided instruments are compared at a common envelope and on more than the knife-edge count. But a policy scored as fully compensating a household on its use-side loss is not thereby compensating that household, and no policy in the scorecard addresses the income channel at all.

There is a countervailing consideration. \citet{pallotti2024} find that euro-area heterogeneity in 2021--23 was driven mainly by nominal net asset positions rather than energy basket composition. If the same holds for the UK, both the use-side channel modelled here and Känzig's employment channel are second to a balance-sheet channel requiring an up-to-date wealth distribution; the model's wealth imputation rests on Wealth and Assets Survey round 7 (2018--20), so this paper cannot speak to it.

\subsection{Three gaps this paper makes visible}

\paragraph{(a) UK wholesale-to-retail pass-through in the price-cap era.} The coefficient $\phi_f$ and lag $\ell$ in equation~\eqref{eq:passthrough} are calibrated from the published structure of the cap rather than estimated, because there is no well-cited econometric estimate of pass-through for the UK under the default tariff cap. The pre-cap literature is not transferable: the cap replaced competitive-retail pricing with a formula, amended repeatedly since 2019, so the estimation problem is regime-switching on a short sample. The data exist --- Ofgem publishes cap decompositions, averaging windows and announcement dates, and forward curves are observable daily --- and identification comes from the announcement calendar itself, the cap being a deterministic function of forwards over a pre-announced window, so the residual is the object of interest for fixed-tariff and out-of-cap customers. This is the most tractable of the three gaps and would tighten every number here.

\paragraph{(b) Energy arrears as an outcome variable.} The outcome here is a cost increase; the observable manifestation of the 2026 shock is debt, with JRF reporting around 1.5 million families in arrears and 3.8 million households owing roughly \pounds 7.5 billion. Arrears depend on liquidity, credit access, supplier forbearance and the prepayment split, so modelling them needs an intertemporal margin the model lacks, and modelling the policy response needs debt-related allowances --- which sit inside the cap, so that arrears feed back into everybody's unit rate. That within-consumer cross-subsidy has not been quantified for the UK, and it would change the policy ranking, since a social tariff that prevents arrears has a fiscal return the scorecard does not credit.

\paragraph{(c) Off-gas-grid and rural horizontal losers.} \citet{douenne2020} shows for France that flat recycling is progressive on average while creating large horizontal losers among rural and off-grid poor households. There is no UK counterpart and this paper does not supply one: the microdata contain no primary-heating-fuel indicator and no gas-grid flag, so losses cannot be split on gas-grid status at all, and no imputed stand-in is offered. The gap is live --- roughly four million UK properties are off the gas grid, heating oil and LPG sit outside the price cap, and these households were repeatedly omitted from 2022 support schemes --- and the finding that motor fuel carries a large part of the loss makes it sharper, since rural transport dependence and off-grid heating coincide in the same households.

\subsection{Where this leaves the macro-to-micro link}

The Seventh Carbon Budget established ``macro model plus distributional annex'' as the expected format for UK official analysis, but the annex is typically produced by a separate method on separate data, so the macro and micro numbers are not mutually consistent. What this paper does is small but structural: it fixes a single explicit transmission mechanism as the bridge --- the cap formula for domestic energy, and for road fuel a pass-through coefficient fitted to observed pump outturns --- so that the household results are a deterministic function of the macro scenario and can be regenerated when the scenario changes. The bridge is one-way, covers only the price channel, and sits alongside the tax-benefit model rather than inside it, because the model has no price channel to host it. Moving it inside, extending it to the income channel, and closing the loop so that the household response feeds back into aggregate demand is the natural continuation, and gap (a) is its prerequisite.
